% !PRINTED 229
\weberpage{229}
Bildet man ausserdem noch $\varphi\psi$, $\varphi^{-1}\psi$, $\varphi\chi\psi$, $\chi\varphi^{-1}\psi$, so folgt:
\[
\begin{aligned}
\varphi\psi &= \Theta\chi\psi\Theta^{-2}, \quad \varphi^{-1}\psi = \Theta^{-2}\chi\psi\Theta, \quad \varphi\chi\psi = \Theta^{-1}\chi\psi\Theta^2, \\
&\chi\varphi^{-1}\psi = \Theta^2\chi\psi\Theta^{-1},
\end{aligned}
\]
woraus man sieht, dass man zu keinem anderen Resultate kommen würde, wenn man die Substitution dritter Ordnung, von der man ausgeht, in der zweiten Form $\Theta^r\chi\psi\Theta^s$ annehmen wollte. Die ganze Gruppe $Q$ ist also durch die angenommene Vierergruppe $D_2$ völlig bestimmt, und man erhält sie in der Form
\weq{3}{
\begin{array}{ll}
1 & \psi \\
\varphi=\Theta\chi\Theta^2, & \varphi\psi=\Theta\chi\psi\Theta^{-2}, \\
\varphi^{-1}=\Theta^{-2}\chi\Theta^{-1}, & \varphi^{-1}\psi=\Theta^{-2}\chi\psi\Theta, \\
\chi & \chi\psi \\
\varphi\chi=\Theta^{-1}\chi\Theta^{-2}, & \varphi\chi\psi=\Theta^{-1}\chi\psi\Theta^2, \\
\varphi^{-1}\chi=\Theta^2\chi\psi\Theta^{-1}, & \varphi^{-1}\chi\psi=\Theta^2\chi\Theta.
\end{array}
}

Man kann diese Gruppe aus den Substitutionen $\varphi, \psi, \chi$ als den Erzeugenden ableiten und sie in die Form setzen:
\weq{4}{Q = \varphi^r\chi^s\psi^t, \qquad \begin{gathered}r=0,1,2,\\ s=0,1; \quad t=0,1.\end{gathered}}

Dass dadurch wirklich eine Tetra\"edergruppe dargestellt ist, ergiebt sich aus den Zusammensetzungen, die man mittelst \S.\,58, (23), (25) leicht aus (1) findet:
\weq{5}{
\begin{aligned}
\psi\varphi &= \varphi\chi\psi, \quad \chi\varphi = \varphi\psi, \\
\psi\varphi^2 &= \chi\varphi^2\psi, \quad \chi\varphi^2 = \varphi^2\chi\psi. \quad \psi\chi = \chi\psi. % [sic: Weber prints a full stop between the two relations, unlike the commas separating the parallel relations; the printed punctuation is preserved.]
\end{aligned}
}

Da in der Tetra\"edergruppe fünf Vierergruppen $D_2$ enthalten sind, so hat die Ikosa\"edergruppe fünf Tetra\"edergruppen zu Theilern. % [sic: Weber prints Tetraëdergruppe; mathematically Ikosaëdergruppe is intended.]
Diese können wir aus $Q$ durch Transformation mittelst der Potenzen von $\Theta$ ableiten und erhalten sie in der Form
\[
\Theta^{-r}Q\Theta^r \quad r=0,1,2,3,4.
\]

Diese Gruppen haben, ausser der Identität, keine Substitution mit einander gemein. Denn die beiden Gruppen $Q$ und $\Theta^{-1}Q\Theta$ haben ausser der Identität nur die beiden Elemente
\[
\Theta^{-2}\chi\Theta^{-1}, \quad \Theta\chi\Theta^2
\]
mit einander gemein, und von diesen kommt keines in $\Theta Q\Theta^{-1}$ vor.

Daraus ergiebt sich nun nach dem Satze 2. in \S.\,28, dass die Ikosa\"edergruppe isomorph ist mit einer Permutationsgruppe
%CONT>>
